% !TEX root = cha-approach.tex
\documentclass[../../main.tex]{subfiles}

\begin{document}

\section{Layout via MDS}

Before generating the paths, the classes have to be placed on the grid. This process consists of four steps:

\begin{enumerate}
    \item Determining the size of each class based on its number of connections. 
    \item Computing a distance matrix between all pairs of classes.
    \item Computing target 2D positions using Multidimensional Scaling (MDS).
    \item Placing the classes on the grid while respecting the placement constraints.
\end{enumerate}

\subsection{Multi-cell Classes}

A class can have many incoming and outgoing connections. Therefore, a single cell may not provide enough space for all its paths. To solve this problem, classes can occupy multiple cells.

Each class is represented as a black square of cells, with side length $S = \lceil k/4 \rceil$, where 
$k$ is the total number of connections of the class, calculated as the sum of its incoming and outgoing connections. For example, a class with up to four connections occupies one cell ($1 \times 1$), while a class with five to eight connections occupies a $2 \times 2$ region and so on. This provides enough space for its paths to enter and leave.

\subsection{Distance Matrix}

To determine the position of the classes, we build an $n \times n$ distance matrix, where $n$ is the total number of classes. 

The distance between Class A and Class B is calculated as the reciprocal of the sum of the weights of their connections:

$$\text{Distance}(A, B) = \frac{1}{\sum \text{Weight of connections between } A \text{ and } B}$$

Different relationship types can be assigned different weights. For example, a method call can have a weight of $2$, while an attribute access can have a weight of $1$. If Class A calls a method in Class B and Class B reads an attribute of Class A, their total connection weight is $2 + 1 = 3$, resulting in a distance of $\frac{1}{3}$.

If two classes have no direct connection, their initial distance is set to a maximum default value.

After filling the matrix with direct distances, we apply the Floyd-Warshall algorithm~\citebibliography{Floy1962a}. This computes the shortest distance between all pairs of classes, enriching the information given to MDS with indirect relationships.

Note that classes that are not connected, either directly or indirectly, preserve the maximum default distance from all other classes.

\subsection{Target Position On Grid}

We apply metric Multidimensional Scaling (mMDS)~\citebibliography{Torg1952a,Doug1998a} to convert the distance matrix into 2D coordinates, in order to preserve the distances between classes as much as possible.

The resulting coordinates are then normalized and scaled to fit the current grid size. An example of assigned positions is presented in \figref{fig:MDS-positions}. The number inside each square indicates how many classes were assigned to that specific position by mMDS, as strongly connected entities may occasionally be mapped to identical coordinates.

\begin{figure}[ht]
    \centering
    \includegraphics[width=0.6\linewidth]{figures/approach/MDS-positions}
    \caption{Example of target positions calculated by mMDS (normalized and scaled to the grid).}
    \label{fig:MDS-positions}
\end{figure}

\subsection{Class Placement Strategy}

After mMDS provides the target positions, classes are placed on the grid one at a time.

\paragraph{Class Placement Order}

The order in which classes are placed may affect the final layout. Our approach supports five ordering strategies:

\begin{itemize}
    \item \textbf{Random}: Classes are placed in a random order.
    \item \textbf{Minimum Connections First}: Classes are sorted by their total number of connections, from the fewest to the most.
    \item \textbf{Maximum Connections First}: Classes are sorted by their total number of connections, from the most to the fewest.
    \item \textbf{BFS Minimum}: Breadth-First Search (BFS) starting from the class with the fewest number of connections.
    \item \textbf{BFS Maximum}: Breadth-First Search (BFS) starting from the class with the most number of connections.
\end{itemize}

To leave enough space for paths, a minimum distance is kept between classes. This parameter defines the minimum number of empty cells required between two classes.

For example, if the minimum distance is $3$, at least three empty cells must separate the borders of two classes horizontally, vertically, and diagonally. 
\SL{C'é bisogno di un'immagine esplicativa?}

\paragraph{Candidates Search}

For each class, the algorithm searches for a valid position as follows:

\begin{enumerate}
    \item The algorithm first tries to place the class at the target position computed by mMDS.
    \item If the target position is not valid, the algorithm searches for a valid position in expanding rings around it, starting with the first ring around the target position.
    \item If one or more valid positions are found in the current ring, the position with the smallest total Manhattan distance to the already placed connected classes is selected.
    \item If no valid position is found in the current ring, the algorithm continues with the next ring.
    \item This process is repeated until at least one valid position is found.
\end{enumerate}

If no valid position is found for a class, the grid is expanded and the placement process restarts.

\end{document}
